\chapter{Data acquisition}
\label{impl:acquisition}
\emph{Mirrors Chapter~\ref{ch:acquisition}.}

\section{Pipeline and tooling}
\label{impl:acqpipeline}
Package \texttt{soaring.acquisition.ffvl}, exposed as two CLI entry points,
\texttt{soaring-para} and \texttt{soaring-delta}, sharing subcommands: \texttt{fetch-xml}
(archive season XMLs), \texttt{download} (fetch IGC files), \texttt{build-catalog}
(rebuild catalog/index), \texttt{status} (per-season coverage), \texttt{verify}
(post-hoc integrity scan, Sec.~\ref{impl:acqintegrity}).

\paragraph{Resumable downloads.} Files already present on disk are skipped, so an
interrupted job resumes freely without re-fetching.

\paragraph{Atomic writes.} Each file is written to a temporary \texttt{.part} and then
renamed, so a file under its final name is always complete -- important on the exFAT
external disk, which has no journaling.

\paragraph{Content validation.} A response is accepted only if it looks like a real IGC
file (an \texttt{A} header record followed by at least one \texttt{B} fix record); HTML
error pages and truncated responses are rejected and retried.

\paragraph{Courteous networking.} A small bounded thread pool, per-request timeouts,
exponential backoff, and a minimum delay between requests avoid hammering the server; a
browser TLS fingerprint is impersonated to pass the site's Cloudflare challenge.

\section{Reproducibility}
\label{impl:acqrepro}
Destination directory (\texttt{data\_root}) set per machine through an environment
variable (\texttt{SOARING\_PARA\_DATA\_ROOT}/\texttt{SOARING\_DELTA\_DATA\_ROOT}) or a
config file shipping a placeholder (\texttt{configs/para\_download.yaml}/
\texttt{configs/delta\_download.yaml}); every command validates the setting at start-up.

\section{Integrity}
\label{impl:acqintegrity}
The \texttt{verify} subcommand re-scans every downloaded file on disk, reusing the same
IGC validation applied at download time, and flags truncated or invalid files, leftover
\texttt{.part} fragments, and read errors.
